\styledchapter{Social Choice with Indifferences, Constrained Allocation}

In Gibbard--Satterthwaite, we assumed everyone had strict preferences, but now we can suppose that agents can be indifferent between allocations, e.g., if they are indifferent between allocations which assign them the same objects.

\section{Introduction and the Equivalence of Constrained Allocation}

We introduce seemingly different problems that will turn out to be equivalent.

\begin{definition}[Social choice problem with indifferences]
	For each agent $i\in N$, suppose we have a partition $A^i \coloneqq \left\{A_1^i, \ldots,A_{m_i}^{i}\right\}$ of $A$ such that if $x,y \in A_\ell^{i}$, then $i$ is indifferent between $x$ and $y$.
	Let $P^i$ be the set of strict preferences on $A^i$.
	Then any $\succ_i \in P^i$ induces a weak ranking $\succsim_i^\star$ on $A$, defined by the following relation for $x \in A_\ell^i$, $y \in A_k^{i}$: \[
		\begin{cases}
			x \succ_i^\star y, \quad A_\ell^{i} \ne A_k^{i} \text{ and } A_\ell^{i} \succ_i A_k^{i} \\
			x \sim_i^\star y, \quad A_\ell^{i}=A_{k}^{i}
		\end{cases}
	.\] 

	A social choice function with indifferences is a function $f: \prod_{i \in N}^{} P^i \to A$ and is defined by \[
		\left( N, A, \left( A^i \right)_{i \in N} \right)
	,\] 
	and the social choice problem is to select such a function given the indifference partitions.
\end{definition}

\begin{definition}[Constrained allocation problem]
	Suppose for each agent $i \in N$, we have a set $X_i$ of available objects. 
	An allocation is a function $\mu: N \to \cup_{i \in N} X_i$ such that for every $i \in N$, $\mu(i) \in X_i$. An allocation specifies for each agent $i$ a feasible object to be received.
	Denote $\mathscr{A} \coloneqq \left\{\mu: N \to \cup_{i \in N}\right\}$ to be the set of allocations.

	A constraint is a set $C \subseteq \mathscr{A}$ that is nonempty.\boxednote{The constraint set is given to us in the allocation problem.} Allocations in $C$ are called \emph{feasible}, while allocations not in $C$ are \emph{infeasible}.
	Let $\tilde{P}_i$ be the set of strict preferences on $X_i$. A mechanism is a function $f: \prod_{i \in N}^{} \tilde{P}_i \to C$.

	The constrained allocation problem is defined by \[
		\left( N, \left( X_i \right)_{i \in N}, C \right)
	.\] 
\end{definition}

\begin{proposition}[Equivalence of social choice with indifferences and constrained allocation]
	Every constrained allocation problem can be written as a social choice problem with indifferences. Conversely, every social choice problem with indifferences can be written as a constrained allocation problem. Under these translations, mechanisms are the same up to outcomes that all agents regard as indifferent.
\end{proposition}
\begin{proof}
	($\implies$) First start with a constrained allocation problem defined on \[
		\left( N, \left( X_i \right)_{i \in N}, C \right)
	.\] 
	Replace each $X_i$ with $\left\{x \in X_i: \mu(i) = x \text{ for some } \mu \in C\right\}$, so every remaining available object is assigned to agent $i$ in some feasible allocation. Then set $A \coloneqq C$.
	For each $i \in N$, define a partition of $A$ by \[
		A^i \coloneqq \left\{ \left\{\mu \in C: \mu(i) = x\right\}: x \in X_i \right\}
	.\] 
	Thus two feasible allocations are in the same indifference class for agent $i$ exactly when they give agent $i$ the same object.\footnote{This is the right indifference class because in the constrained allocation problem, agent $i$ only has a strict preference over her own set $X_i$ of objects and does not care about others' objects.}

	A strict preference $\succ_i$ on $X_i$ induces a strict preference $\succ_i^\star$ on $A^i$ by \[
		\left\{\mu \in C: \mu(i) = x\right\} \succ_i^\star \left\{\mu \in C: \mu(i) = y\right\} \iff x \succ_i y
	.\] 
	This gives a canonical identification of the constrained allocation preference domain $\tilde{P}_i$ with the social choice preference domain $P^i$.
	Hence a constrained allocation mechanism $f: \prod_{i \in N}^{} \tilde{P}_i \to C$ can be regarded as the social choice function with indifferences $f^\star: \prod_{i \in N}^{} P^i \to A$ given by \[
		f^\star\left( \succ_1^\star,\ldots,\succ_{\left| N \right|}^\star \right)
		=
		f\left( \succ_1,\ldots,\succ_{\left| N \right|} \right)
	.\] 
	In this direction, the equivalence of mechanisms is immediate because the social outcomes are exactly the feasible allocations: $A=C$.

	\medskip

	($\impliedby$) Conversely, start with a social choice problem with indifferences defined on $\left( N,A,\left( A^i \right)_{i \in N} \right)$.
	Set $X_i \coloneqq A^i$, so agent $i$'s available objects are her indifference classes.
	With this definition, the constrained allocation preference domain $\tilde{P}_i$ is exactly $P^i$, the set of strict preferences over $A^i$.
	Note that any allocation $\mu$ chooses one class $\mu(i) \in A^i$ for each agent $i$.
	Define the constraint as a subset of all possible allocations $\mathscr{A} \coloneqq \left\{\mu: N \to \times_{i \in N} X_i \right\}= \left\{\mu: N \to \times_{i \in N} A_i\right\}$ by \[
		C \coloneqq \left\{\mu \in \mathscr{A}: \bigcap_{i \in N}^{} \mu(i) \ne \O \right\}
	.\] 
	A tuple $(\mu(i))_i	\in \prod_{i \in N}^{} A^i$ of assigned indifference classes is feasible if there exists at least one original outcome $a \in A$ such that $a$ is in every assigned indifference class.

	Thus any social choice function $f: \prod_{i \in N}^{} P^i \to A$ induces a constrained allocation mechanism $\hat{f}: \prod_{i \in N}^{} \tilde{P}_i \to C$ by \[
		\hat{f}(\succ)(i) = A_\ell^i
		\quad \text{where } f(\succ) \in A_\ell^i
	.\] 
	$\hat{f}(\succ)(i)$ is the indifference class $A_\ell^{i}$ in $i$'s partition of $A$, which is agent $i$'s object in the constructed constrained allocation problem.

	To show that no preference-relevant information is lost in this translation, take a constrained allocation mechanism $g: \prod_{i \in N}^{} \tilde{P}_i \to C$ in the constructed allocation problem.
	Since each feasible allocation $g(\succ)$ is a tuple of indifference classes with nonempty intersection, choose for each $\mu \in C$ some representative $r(\mu) \in \cap_{i \in N} \mu(i)$ and define the social choice function $\hat{g}: \prod_{i \in N}^{} P^i \to A$ by \[
		\hat{g}(\succ) \coloneqq r(g(\succ))
	.\] 
	The choice of representative does not matter for preferences, since every representative lies in the same indifference class for every agent.
\end{proof}

\begin{example}[Examples of Equivalence]
	Consider:
	\begin{enumerate}[label = (\roman*)]
		\item House allocation: Suppose we have a set of possible houses $H$, and put $X_i \coloneqq H$ for every $i$. The constrained set of allocations is \[
			C = \left\{\mu: N \to H: \mu(i) \ne \mu(j), \text{ if } i \ne j\right\}
		.\] 
		\item School choice: Denote the set of schools $X$, and put $X_i \coloneqq X$ for all $i$. Each school $x \in X$ has a capacity $q_x \ge 0$, and the constrained set is \[
			C = \left\{\mu: N \to X: \left| \mu^{-1}(x) \right| \le q_x, \forall x\right\}
		.\] 
		\item One-sided matching:\footnote{More commonly called the roommates problem.} Suppose $X_i \coloneqq N$ for every $i$ and we have the constrained set \[
			C = \left\{\mu: N \to N : \mu\circ \mu (i) = i, \forall i\right\}
		.\] 
		$C$ requires that if $i$ is matched with $j$, $j$ must be matched with $i$. But it allows for people to be matched with themselves, i.e. unmatched.
		\item Two-sided matching: Consider $N = M \amalg W$. Set \[
			X_i = \begin{cases}
				W, \quad i \in M \\ M, \quad i \in W
			\end{cases}
		,\] and put \[
			C = \left\{\mu: N \to N: \mu \circ \mu(i) = i \text{ and } \mu(i) \in M \text{ if } i \in W, \forall i\right\} 
		.\] 
		\item Social choice: Put $X_i \coloneqq A$ for every $i$ and let \[
			C = \left\{\mu: N \to A : \mu(i) = \mu(j), \forall i,j\right\}
		.\] Thus the social choice problem we have been thinking of this whole time is a special case of the constrained allocation problem.
	\end{enumerate}
\end{example}

We reintroduce two concepts in this new setting.

\begin{definition}[Strategy-proof, Pareto efficient]
	In the setting of a constrained allocation problem, a mechanism $f: \prod_{i \in N}^{} \tilde{P}_i \to C$ is strategy-proof if there does not exist $i, \succ_i, \succ_i', \succ_{-i}$ such that \[
		f\left( \succ_i', \succ_{-i} \right)(i) \succ_i f\left( \succ_i, \succ_{-i} \right)(i)
	.\] 

	$f$ is Pareto efficient if for all $\succ \in \prod_{i \in N}^{} \tilde{P}_i$, there does not exist $\mu \in C$ such that:
	\begin{enumerate}[label = (\roman*)]
		\item For every $j \in N$, $\mu(j) \succsim_j f(\succ)(j)$ and 
		\item For some $i \in N$, $\mu(i) \succ_i f(\succ)(i)$.
	\end{enumerate}
\end{definition}

\section{Two Agent Constrained Allocation}

Now, assume $N = \left\{1,2\right\}$.
Fix the possible sets $X_1,X_2$ and a constraint $C$.

\begin{definition}[Local dictator]
	A local dictator at an infeasible allocation $(x,y) \not\in C$ is an agent who gets to keep her top choice when $x,y$ are top-ranked by agents 1 and 2 respectively.
\end{definition}

We will describe all strategy-proof and Pareto efficient mechanisms for $C$.

\begin{proposition}[Two-agent local dictatorships] \label{4-13-a}
	Let $G$ be the graph whose nodes are infeasible allocations and whose edges connect two infeasible allocations that share a coordinate.
	After removing objects that can never be assigned to an agent in any feasible allocation, a mechanism $f$ is strategy-proof and Pareto efficient if and only if it is a local dictatorship: each connected component of $G$ is assigned a local dictator, and for every profile with top-ranked objects $(x,y)$, the mechanism gives both agents their top choices if $(x,y) \in C$; otherwise, the local dictator for the component containing $(x,y)$ keeps that agent's top choice and the other agent receives her favorite object compatible with the dictator's top choice.
\end{proposition}

\begin{lemma} \label{4-13-1}
	Suppose that there exists $x \in X_i$ for some $i$ such that $\mu(i) \ne x$ for every $\mu \in C$. 
	If $f$ is strategy-proof and Pareto efficient, then for $\succ_i$ and $\succ_i'$ that differ only in the position of $x$, then \[
		f(\succ_i, \succ_j)(i) = f\left( \succ_i', \succ_j \right)(i)
	.\] 
	Furthermore, $j$'s object is unchanged by $i$'s misreport: \[
		f\left( \succ_i,\succ_j \right)(j)
		= 
		f\left( \succ_i',\succ_j \right)(j)
	.\] 
\end{lemma}
\begin{proof}
	If the objects are not the same, then $i$ has a strict preference between the two objects. 
	Suppose \[
		f(\succ_i, \succ_j)(i) \succ_i f\left( \succ_i', \succ_j \right)(i)
	.\] Since neither side is $x$ and $\succ_i,\succ_i'$ differ only in the position of $x$, we also have \[
		f(\succ_i, \succ_j)(i) \succ_i' f\left( \succ_i', \succ_j \right)(i)
	,\] 
	which violates strategy-proofness.

	If instead \[
		f(\succ_i, \succ_j)(i) \prec_i f\left( \succ_i', \succ_j \right)(i)
	,\] 
	strategy-proofness is immediately violated.

	\medskip

	For the second claim, suppose for contradiction (without loss) that \[f\left( \succ_i, \succ_j \right)(j) \succ_j f\left( \succ_i',\succ_j \right)(j).\]
	Then $f(\succ_i', \succ_j)$ is not efficient, as $f(\succ_i, \succ_j) \in C$ and $f(\succ_i,\succ_j) \sim_i f\left( \succ_i', \succ_j \right)$ by the first claim.
\end{proof}

We can write an allocation as a pair $(x,y)$. By Lemma (\ref{4-13-1}), for each $i$, we can get rid of all $x \in X_i$ that can never be chosen\boxednote{This is true because if there does not exist such a $y'$, $\mu(i) \ne x$ for all $\mu$, but we already got rid of all such $x$.} and assume that for every $(x,y) \not\in C$, there exist $x', y'$ such that $(x',y) \in C$ and $\left( x,y' \right) \in C$.

Let $\overline{C} \subseteq C$ be the set of infeasible allocations.

\begin{lemma} (\ref{5-23-1})
	For any strategy-proof and efficient mechanism $f$, each infeasible allocation $\left( x,y \right) \not\in  C$ is associated with a local dictator. In particular, there exists a function $D: \overline{C} \to \left\{1,2\right\}$ that describes who is the local dictator at every infeasible allocation.

	Furthermore, if $(x,y'') \not\in C$ as well, then \[
		D\left( \left( x,y \right) \right) = 1 \implies D\left( \left( x,y'' \right) \right) = 1
	.\] 
\end{lemma}

\begin{proof}	
	Suppose $(x,y)$ is not feasible and there exist $x',y'$ such that $(x',y) \in C$ and $(x,y') \in C$. Consider the preferences
	\begin{align*}
		&\succ_1 : x \succ_1 x' \succ_1 \cdots \\
		&\succ_2: y \succ_2 y' \succ_2 \cdots
	.\end{align*}
	We know that $f(\succ_1,\succ_2) \in \left\{\left( x,y' \right), \left( x',y \right)\right\}$ by Pareto efficiency of $f$.
	Suppose without loss that we side with agent 1 and $f(\succ_1,\succ_2) = (x,y')$.

	Consider another preference $\succ_2'$ where $y$ is top-ranked. We know by strategy-proofness that $f(\succ_1,\succ_2')(2) \ne y$. This implies $f(\succ_1,\succ_2')(1) = x$ by Pareto efficiency, as otherwise, we could have put $f(\succ_1,\succ_2') = (x',y)$, making agent 1 no worse off and agent 2 strictly better off.
	By one more application of strategy-proofness, $f(\succ_1',\succ_2')(1) = x$ for any $\left( \succ_1',\succ_2' \right)$ such that 1 top-ranks $x$ and 2 top-ranks $y$, as if agent 1 does not get $x$ by ranking $x \succ_1' x'' \succ_1' \cdots$, she could deviate to reporting $x \succ_1 x' \succ_1 \cdots$ and get $x$, which is her top choice under $\succ_1'$. By Pareto efficiency, this implies that agent 2 gets his favorite object compatible with 1 getting $x$.

	We have shown that agent $1$ is a local dictator over $(x,y)$.
	We can repeat this procedure for every infeasible allocation,
	put \[
		D: \overline{C} \to \left\{1,2\right\}
	,\] 
	where the function indicates who we sided with in getting his or her top-ranked object.

	\medskip

	For the second claim, suppose not, so $D\left( \left( x,y'' \right) \right) = 2$. Consider the preferences 
	\begin{align*}
		&\succ_1: x \succ_1 \cdots \\
		&\succ_2: y \succ_2 y'' \succ_2 \cdots \\
		&\succ_2': y'' \succ_2' \cdots
	.\end{align*}
	Under $(\succ_1,\succ_2)$ then 2 gets something worse than $y''$, since 1 is the local dictator at $(x,y)$ and $(x,y'')$ is infeasible. But 2 can misreport $\succ_2'$ to get $y''$, as he is the dictator at $(x,y'')$, which violates strategy-proofness.

	The same argument works for infeasible allocations that share the second coordinate instead of the first coordinate. Repeating this argument along a connected path shows that all infeasible allocations in the same connected component of $G$ have the same local dictator.
\end{proof}

\propositionproofstyle
\begin{proof}[Proof of Proposition \ref{4-13-a}]
	The first lemma lets us ignore any object that is never assigned to an agent in a feasible allocation, since changing the position of such an object cannot affect either agent's allocation.
	Thus, for any infeasible $(x,y)$, there exist $x'$ and $y'$ such that $(x',y) \in C$ and $(x,y') \in C$.

	The second lemma then assigns a local dictator to every infeasible allocation and shows that this assignment is constant across infeasible allocations that share a coordinate. Equivalently, it is constant on every horizontal or vertical edge of $G$, and therefore on every connected component of $G$.
	Hence any strategy-proof and Pareto efficient mechanism must be a local dictatorship.

	Conversely, suppose $f$ is a local dictatorship.
	It is Pareto efficient because if the top-ranked pair $(x,y)$ is feasible, both agents receive their top choices; if $(x,y)$ is infeasible and agent 1 is the local dictator, then 1 receives $x$ and agent 2 receives her favorite object among the feasible allocations with first coordinate $x$, so no feasible allocation can improve 2 without making 1 worse off. The case where agent 2 is the local dictator is symmetric.

	To see strategy-proofness, first suppose the reported top-ranked objects are feasible. Then each truthful agent receives that agent's favorite object and cannot gain by deviating. If the reported top-ranked pair is infeasible, the local dictator cannot gain because the dictator keeps the dictator's favorite object. The non-dictator receives the best object compatible with the dictator's reported top choice. 
	A deviation by the non-dictator to another infeasible allocation means the dictator is the same by Lemma (\ref{5-23-1}), so the non-dictator cannot improve further on his or her best compatible object. 
	A deviation that makes the top-ranked pair feasible gives the dictator the dictator's top choice and therefore gives the deviating non-dictator an object compatible with that same top choice, again no better than what the non-dictator already received. Thus no unilateral deviation is profitable.
\end{proof}
